\documentclass[journal]{IEEEtran}

\usepackage{amsmath}
\usepackage{amssymb}
\usepackage{amsthm}
\usepackage{algorithm}
\usepackage{algorithmic}
\usepackage{booktabs}
\usepackage{graphicx}
\usepackage{url}

\newtheorem{theorem}{Theorem}
\newtheorem{lemma}{Lemma}
\newtheorem{proposition}{Proposition}
\newtheorem{corollary}{Corollary}
\newtheorem{conjecture}{Conjecture}

\newcommand{\eqdef}{:=}

\begin{document}

\title{Fair, Non-Progressive Influence Control Under Backfire and Parameter Uncertainty}

\author{Quang-Vinh~Dang
\thanks{Q.-V. Dang is with the British University Vietnam, Hung Yen, Vietnam (e-mail: vinh.dq4@buv.edu.vn).}}

\markboth{IEEE Transactions on Computational Social Systems}{}

\maketitle

\begin{abstract}
The classical influence maximization guarantee of Kempe, Kleinberg, and Tardos
rests on four assumptions: progressive diffusion, a monotone submodular
objective, known influence parameters, and no fairness requirement. Real
deployments routinely violate all four at once: influence decays and can be
actively reversed (backfire, or negative word-of-mouth), edge probabilities
are never known exactly, and social interventions must not neglect a
disadvantaged subpopulation. This paper studies the combined relaxation. We
prove that a sandwich approximation built from two monotone-submodular
surrogates yields a $\rho(1-1/e)$ guarantee for \emph{every} backfire
intensity $\beta \in [0,1)$, degenerating gracefully to the tight $1-1/e$
bound at $\beta=0$ and admitting a per-instance, directly computable
certificate. We derive a closed-form Lagrangian index for the resulting
2-state/3-action per-node control problem, replacing a bisection-plus-value-iteration
inner loop with an $O(Km+n\log n)$-per-round procedure that we measure to be
roughly two orders of magnitude faster than the bisection oracle it replaces.
We show that $(1-1/e)$-hardness survives backfire exactly, and we state --
without overclaiming -- a conditional, not-yet-proven result suggesting a
further constant-factor worsening for a plausible but unverified property of
Feige's hard instances. Finally, we identify and correct a population-weighting
pathology in an existing isoelastic-welfare fairness objective that
systematically under-protects small groups, and replace it with an
egalitarian reweighting. We implement all of this in a control policy,
MF-BWI-Fair, that combines mean-field belief propagation, Bayesian parameter
tracking, the closed-form index, and the corrected fairness reweighting.
Preliminary experiments on a synthetic stochastic block model benchmark
show MF-BWI-Fair sustaining spread under backfire and recovery where
one-shot baselines -- including independently reimplemented published
algorithms (IMM, He--Kempe's robust Saturate Greedy) -- degrade sharply, and
show the egalitarian fix raising worst-off-group reach at negligible spread
cost. A full multi-graph validation and a benchmark against the original
IMM implementation are in progress and are not reported here.
\end{abstract}

\begin{IEEEkeywords}
Influence maximization, submodular optimization, restless bandits, fairness,
robust optimization, social networks.
\end{IEEEkeywords}

\section{Introduction}

Kempe, Kleinberg, and Tardos~\cite{kempe2003maximizing} showed that under the
Independent Cascade (IC) and Linear Threshold (LT) models, the expected spread
of a seed set is monotone and submodular, so the greedy algorithm achieves a
$(1-1/e)$-approximation via the general result of Nemhauser, Wolsey, and
Fisher~\cite{nemhauser1978analysis}, and this ratio is tight under standard
complexity assumptions~\cite{feige1998threshold}. That guarantee, and the
large body of scalable algorithms it spawned -- CELF's lazy-forward
speedup~\cite{leskovec2007cost} and the near-linear-time RIS/martingale
algorithm IMM~\cite{tang2015influence,borgs2014maximizing} chief among them --
rests on four assumptions taken together: (A1) diffusion is
\emph{progressive} (once active, always active); (A2) the spread function is
\emph{monotone submodular}; (A3) influence parameters are \emph{known}
exactly; and (A4) the objective has \emph{no fairness structure} -- every
node's activation counts equally toward one aggregate number.

Each assumption is routinely false in the settings influence maximization is
used to model. Interventions decay: a vaccinated or informed individual can
relapse or forget, and messaging campaigns require sustained reinforcement,
not a one-time seeding. Influence can be actively \emph{negative}: word of
mouth can backfire, and an over-exposed contact can turn an active neighbor
against the message rather than merely fail to convert
it~\cite{chen2011influence,he2012influence,li2014influence}. Edge
probabilities are estimated from noisy or incomplete cascade data, never
known exactly. And in the social-intervention settings that motivate much of
this work in the first place -- HIV-prevention and stable-housing
information campaigns among homeless youth, the application context
Rahmattalabi et al.\ build their fairness-aware
formulation around~\cite{rahmattalabi2021fair} -- treating every node as
interchangeable is not a simplification but a policy failure: a mechanism
that maximizes aggregate reach while systematically under-serving one
subpopulation is not simply "unfair" in the abstract, it directly
undermines the deployment's stated public-health goal.

Prior work has relaxed these assumptions one, or occasionally two, at a
time (Section~\ref{sec:related}). This paper studies all four relaxations
\emph{together}: a model with recovery and backfire-driven deactivation
(relaxing A1 and, as we show, A2), Bayesian tracking of unknown influence
and recovery parameters (relaxing A3), and an egalitarian-welfare
fairness objective over groups (relaxing A4). Our literature search found no
existing paper that combines all four axes; the two closest works each
combine exactly two (Section~\ref{sec:related}).

\subsection*{Contributions}
\begin{itemize}
\item \textbf{A formal layered-semantics model} (Section~\ref{sec:model})
  in which an active node survives a round unless it spontaneously recovers
  or is backfired by an active in-neighbor, and any node -- active or not --
  can be (re-)activated by a successful positive trial from an active
  in-neighbor.
\item \textbf{A proven submodularity result for recovery alone}
  (Lemma~\ref{lem:recovery-submodular}): spontaneous, state-independent
  recovery preserves monotone submodularity of cumulative spread on any
  graph, via a layered live-edge-graph reachability argument. This isolates
  recovery as the "free" relaxation; neighbor-dependent backfire is what
  destroys the coverage structure.
\item \textbf{A sandwich approximation theorem} (Theorem~\ref{thm:sandwich},
  \emph{proven}, instantiating the general sandwich technique
  of~\cite{lu2015competition} together with~\cite{nemhauser1978analysis}):
  for \emph{every} backfire intensity $\beta \in [0,1)$, a greedy solution
  taken over two monotone-submodular surrogates and the raw objective
  achieves $\rho(1-1/e)\,\mathrm{OPT}_\beta$, where $\rho$ is an
  instance-dependent sandwich ratio that degenerates to the tight $1-1/e$
  bound at $\beta=0$. The guarantee is also directly instance-certifiable:
  a returned set ships its own approximation ratio, computable by simulating
  the two $\beta=0$ surrogates.
\item \textbf{A closed-form Lagrangian index} (Propositions~\ref{prop:closedform}--\ref{prop:indexability},
  Theorem~\ref{thm:complexity}, \emph{proven} except for indexability itself,
  which is checked numerically and explicitly flagged as such) for the
  resulting 2-state/3-action per-node control problem, reducing the
  per-round decision from bisection with per-node value iteration to
  $O(Km+n\log n)$ -- an approximately $130\times$ speedup measured in this
  project's implementation.
\item \textbf{Hardness results that preserve honest status labels}:
  Theorem~\ref{thm:hardness} \emph{proves} that $(1-1/e)$-inapproximability
  survives backfire exactly, for every fixed $\beta$; Proposition~\ref{prop:conditional}
  is stated and discussed as a \emph{conditional, unproven} result --
  consistent with its labeling in the underlying project notes -- suggesting
  (not proving) a further constant-factor worsening contingent on one
  plausible but unverified property of Feige's hard-instance construction.
\item \textbf{A corrected fairness mechanism}: we identify a measured
  population-weighting pathology in the isoelastic welfare objective
  of~\cite{rahmattalabi2021fair} (small groups are protected only once
  under-served by a factor proportional to a group-size ratio) and replace
  it with an egalitarian (population-unweighted) reweighting that protects
  the worst-off group at a neutral default setting.
\item \textbf{MF-BWI-Fair}, a control policy combining mean-field belief
  propagation, Bayesian posterior tracking, the closed-form Lagrangian
  index, and the corrected welfare reweighting, together with a preliminary
  comparison against six algorithms including independently-verified
  reimplementations of two published baselines.
\end{itemize}

\section{Related Work}
\label{sec:related}

\subsection{Classical influence maximization}
Kempe, Kleinberg, and Tardos~\cite{kempe2003maximizing} established
monotone submodularity of expected spread under IC/LT and the resulting
$(1-1/e)$ greedy guarantee~\cite{nemhauser1978analysis}, tight under
$\mathrm{P}\ne\mathrm{NP}$~\cite{feige1998threshold}. CELF~\cite{leskovec2007cost}
gives a lazy-forward speedup exploiting the same submodularity. IMM~\cite{tang2015influence},
building on the RIS/reverse-reachable-set duality of~\cite{borgs2014maximizing},
gives an $O((k+\ell)(n+m)\log n/\varepsilon^2)$ near-linear-time
$(1-1/e-\varepsilon)$-approximation and is, to our knowledge, the most-cited
scalable classical algorithm with publicly available original
code~\cite{chen2018issue}.

\subsection{Fair / group-fairness-constrained influence maximization}
Tsang et al.~\cite{tsang2019group} introduce maximin and diversity-constraint
fairness notions as a multi-objective submodular problem. Fish et
al.~\cite{fish2019gaps} study a maximin welfare objective directly and prove
it is not monotone submodular in general. Farnadi et
al.~\cite{farnadi2020unifying} unify several fairness notions as exact MILP
formulations at small scale, with no submodular approximation guarantee.
Becker et al.~\cite{becker2021fairness} show deterministic maximin-fair
seeding is hard to approximate but that randomized seeding recovers a
$(1-1/e-\varepsilon)$-type guarantee. Rahmattalabi et
al.~\cite{rahmattalabi2021fair} propose an isoelastic (CES) welfare objective
over per-group reach and prove monotone submodularity via the
concave-nondecreasing-composed-with-monotone-submodular lemma
of~\cite{lin2011class}, giving an exact $(1-1/e)$ guarantee -- but only for
the classical progressive IC model; their proof does not address
deactivation or negative influence. This paper adopts their objective
\emph{form}, not their guarantee, and additionally corrects a
population-weighting side effect of their exact formula (Section~\ref{sec:fairness-fix}).
Rui et al.~\cite{rui2023scalable} and Feng et al.~\cite{feng2023influence}
extend fairness to scalable and data-driven/learning-based settings,
respectively; the latter is out of scope for the classical-algorithms focus
here.

\subsection{Non-progressive / recovery-capable influence maximization}
Chan and Ning~\cite{chan2015influence} show that non-progressive dynamics
under the LT model lose submodularity on general graphs, recovering it only
on DAGs. Golnari et al.~\cite{golnari2014revisiting} propose the "Heat
Conduction" model, proving submodularity under a specially engineered
non-progressive rule. Lou et al.~\cite{lou2014modeling} give a continuous-time
two-state Markov estimation model without a submodularity result. Zahoor et
al.~\cite{zahoor2024influence} and Hui et al.~\cite{hui2024nonprogressive}
study reactivation dynamics, the latter via deep reinforcement learning
(out of scope here). Tan{\i}nm{\i}s et al.~\cite{taninmis2019influence} study
an adversarial deactivation game, a different framing from the spontaneous,
non-adversarial backfire studied here. Our Lemma~\ref{lem:recovery-submodular}
gives a stronger/different result than Chan and Ning's DAG-restricted
finding: spontaneous, state-independent recovery alone preserves monotone
submodularity on \emph{any} graph, via layered-DAG reachability rather than
threshold-rule structure.

\subsection{Non-monotone / non-submodular / negative-influence models}
Bharathi et al.~\cite{bharathi2007competitive} give the first competitive
cascade model. Chen et al.~\cite{chen2011influence} extend IC with a quality
factor allowing activated nodes to flip negative, engineered to remain
submodular. He et al.~\cite{he2012influence} introduce a two-competing-cascade
"sandwich approximation" as a scalability device, whose general form is due
to Lu, Chen, and Lakshmanan~\cite{lu2015competition}; our
Theorem~\ref{thm:sandwich} instantiates the latter's general
monotone-submodular bracket $L\le f\le U$, applied to a materially different
bracket -- a recovery-only process, not a two-cascade one. Li et
al.~\cite{li2014influence} give a signed-network model (IC-P) that also
remains submodular by construction. Buchbinder et
al.~\cite{buchbinder2015tight} and Tukan et al.~\cite{tukan2024practical} give
tight and practical constant-factor guarantees for unconstrained/cardinality-constrained
\emph{non-monotone submodular} maximization -- a different relaxation
axis (non-monotonicity of a still-submodular function) from ours. Bian et
al.~\cite{bian2017guarantees} give a submodularity-ratio/curvature-based
greedy guarantee for non-submodular functions, superseded here by the
sandwich approach, which holds for all $\beta$ rather than requiring an
estimated ratio. Horel and Singer~\cite{horel2016maximization} give an
exponential query-complexity lower bound for functions
$\varepsilon$-far from submodular once $\varepsilon=\omega(n^{-1/2})$, but
only for an \emph{adversarial} perturbation; our backfire perturbation is
structured (a bounded extra kill rate), which is why Theorem~\ref{thm:sandwich}
escapes this limitation (Section~\ref{sec:hardness}).

\subsection{Bayesian, robust, and online influence maximization under uncertainty}
He and Kempe~\cite{he2016robust} formalize worst-case-across-scenarios robust
influence maximization and give a bicriteria Saturate Greedy algorithm with a
provable budget-blowup/quality trade-off; no official code is available for
this specific paper. The dominant paradigm for online uncertainty is
combinatorial-bandit learning~\cite{vaswani2015influence,wen2017online,lei2015online,vaswani2017modelindependent,chen2013combinatorial,chen2016combinatorial}
or one-shot distributionally robust
optimization~\cite{chen2020correlation,staib2019distributionally}. A
genuinely single-shot Bayesian formulation -- integrating expected spread
over a prior without adaptivity or worst-case hedging -- exists only
embedded inside bandit posterior-update machinery, not as a standalone
anchor result; this is the gap our Bayesian tracker sits in.

\subsection{Restless-bandit sequential control}
Mate et al.~\cite{mate2020collapsing} give binary-action Whittle-index
results for collapsing bandits, explicitly leaving multi-action indexability
future work. Killian, Perrault, and Tambe~\cite{killian2021beyond} give the
Lagrangian-relaxation template for general multi-action restless bandits that
our closed-form index specializes and simplifies for the specific 2-state/3-action
case, avoiding the general indexability problem they call "notoriously
difficult." Ou et al.~\cite{ou2022networked} propose an alternative to
mean-field network coupling that we verified does not apply here: their
coupling is through a deterministic, known action-vector, not through
another arm's actual random state, which is the genuinely bilinear coupling
our setting requires. Li and Varakantham~\cite{li2022efficient} study
fairness constraints in independent-arm restless bandits, confirming
fairness-in-RMAB as an active but separate line of work.

\subsection{Positioning}
No paper we found combines group fairness, non-progressive/recovery
dynamics, non-monotone backfire influence, and Bayesian parameter
uncertainty. The two closest works each combine exactly two of these four
axes, both very recent: a 2026 PACMMOD paper combining robust fairness with
multiple-community-partition (parameter) uncertainty~\cite{robustfair2026pacmmod},
and Fang et al.'s 2026 scalable fair influence \emph{blocking} maximization,
combining fairness with non-monotone/blocking-style negative
influence~\cite{fang2026scalable}. Separately, no paper we found connects
submodular/RIS-style influence-maximization theory to restless-bandit
sequential control theory at all; Section~\ref{sec:index} is, to the depth
of our search, the first to do so for a node-level activation/maintenance
control problem.

\section{Problem Formulation}
\label{sec:model}

Let $G=(V,E)$ be a directed graph, $n=|V|$, $m=|E|$, with a partition of $V$
into groups $G_1,\dots,G_M$. Rounds are $t=0,1,\dots,T$; the active set
$A_t\subseteq V$ starts at $A_0=S$, the seed set. Each round draws
independent randomness: for every edge $(u,v)$, a positive trial
$X^t_{uv}\sim\mathrm{Bern}(p_{uv})$ and a backfire trial
$Y^t_{uv}\sim\mathrm{Bern}(\beta p_{uv})$, for a global backfire intensity
$\beta\in[0,1)$; for every node $v$, a recovery trial
$R^t_v\sim\mathrm{Bern}(q_v)$.

\textbf{Transition (layered semantics).} $v\in A_{t+1}$ iff
\begin{equation}
\begin{split}
\big[\,v\in A_t \wedge R^t_v=0 \wedge \forall u\in A_t\cap N_{\mathrm{in}}(v): Y^t_{uv}=0\,\big]\\
\vee\ \big[\,\exists u\in A_t\cap N_{\mathrm{in}}(v): X^t_{uv}=1\,\big].
\end{split}
\end{equation}
An active node stays active unless it recovers or is backfired by an active
in-neighbor, and any node -- active or not -- is (re-)activated by a
successful positive trial from an active in-neighbor; in particular a
positive trial can \emph{rescue} an active node in the same round it would
otherwise have recovered or been backfired. This rescue clause is exactly
what makes the one-step map $A_t\mapsto A_{t+1}$ monotone in $A_t$ at
$\beta=0$, which the submodularity result of Section~\ref{sec:theory}
rests on; for $\beta>0$ the map is not monotone, since the backfire clause
is anti-monotone.

The objective is cumulative expected spread,
\begin{equation}
f_\beta(S) \;=\; \mathbb{E}\Big[\sum_{t=1}^T |A_t|\Big], \qquad A_0=S,
\end{equation}
time-averaged spread being the same quantity up to the constant $1/T$.

\textbf{Fairness.} Let $u_g$ denote group $g$'s time-averaged realized reach
fraction (active fraction of $G_g$, averaged over rounds). Following an
isoelastic (constant elasticity of substitution) welfare
form~\cite{rahmattalabi2021fair}, define
\begin{equation}
W_\alpha(u) = \sum_g N_g\,\frac{u_g^\alpha}{\alpha}\ (\alpha\ne0), \qquad
W_0(u) = \sum_g N_g \log u_g,
\end{equation}
population-weighted by group size $N_g$; $\alpha\to1$ is utilitarian,
$\alpha=0$ is proportional fairness, and $\alpha\to-\infty$ approaches
leximin. As detailed in Section~\ref{sec:fairness-fix}, this project
identifies and corrects a measured pathology in the population-weighted
form and adopts, by default, the egalitarian (population-unweighted)
variant $W_\alpha(u)=\sum_g u_g^\alpha/\alpha$ instead, which counts each
group once regardless of size.

\section{Theoretical Results}
\label{sec:theory}

Statements below preserve the status labels of the underlying project
notes exactly: \emph{proven} (theorem, proposition, lemma, with a proof or
complete proof sketch), \emph{cited} (an external result relied upon), or
\emph{conjectural/conditional} (stated precisely, with the reason it is
not yet proven). No conjectural or conditional statement is used as a
premise for any theorem.

\subsection{Submodularity of the recovery-only process}

\begin{lemma}[Layered live-edge representation]
\label{lem:layered}
Fix $\beta=0$. Build a layered graph $\mathcal{L}$ on $V\times\{0,\dots,T\}$
with, for each round $t$, a self-edge $(v,t)\to(v,t{+}1)$ live iff
$R^t_v=0$, and a cross-edge $(u,t)\to(v,t{+}1)$ for each $(u,v)\in E$, live
iff $X^t_{uv}=1$. Then for every realization of the randomness,
$A_t = \{v : (v,t) \text{ reachable from } S\times\{0\}\text{ along live edges}\}$.
\end{lemma}
\emph{Proof.} By induction on $t$; at $t=0$ both sides equal $S$, and the
inductive step is exactly the layered-semantics transition rule read as a
reachability statement. $\blacksquare$

\begin{lemma}[Recovery alone preserves monotone submodularity]
\label{lem:recovery-submodular}
For $\beta=0$ and any recovery rates $q_v\in[0,1]$, $f_0(S)$ is monotone and
submodular in $S$.
\end{lemma}
\emph{Proof.} For each fixed realization, $|A_t|$ is the number of
layer-$t$ nodes reachable from $S\times\{0\}$ in a fixed DAG (Lemma~\ref{lem:layered}) --
a coverage function of $S$, hence monotone submodular; sums over $t$ and
expectations over realizations of monotone submodular functions remain
monotone submodular. $\blacksquare$

This isolates exactly which relaxation is "free": spontaneous,
state-independent recovery is harmless (it is a random self-edge and
reachability stays a coverage function), while neighbor-dependent
deactivation (backfire) destroys the coverage structure, since an active
in-neighbor can now \emph{remove} a node from the reachable set. This
sharpens the picture in Chan and Ning~\cite{chan2015influence}, who lose
submodularity under non-progressive \emph{threshold} rules where
deactivation depends on neighbor states; recovery alone needs no such
structure to stay submodular, on any graph.

\subsection{A sandwich approximation for all backfire intensities}

Let $U(S)\eqdef f_0(S)$ (the $\beta=0$ process at the true recovery rates).
For node $v$, let
$\delta_v \eqdef 1-\prod_{u\in N_{\mathrm{in}}(v)}(1-\beta p_{uv}) \le \beta\Delta_{\mathrm{in}}$,
where $\Delta_{\mathrm{in}}\eqdef\max_v\sum_{u\in N_{\mathrm{in}}(v)}p_{uv}$,
and let $L(S)$ be the $\beta=0$ process under inflated recovery
$q'_v\eqdef1-(1-q_v)(1-\delta_v)$.

\begin{lemma}[Upper coupling]
\label{lem:upper}
Under shared randomness, $A^\beta_t\subseteq A^0_t$ for all $t$; hence
$f_\beta(S)\le U(S)$.
\end{lemma}

\begin{lemma}[Lower coupling]
\label{lem:lower}
There is a coupling under which $A^L_t\subseteq A^\beta_t$ for all $t$;
hence $L(S)\le f_\beta(S)$.
\end{lemma}

Both couplings are proven by induction, replacing the per-edge backfire
trials with a single uniform draw compared against the aggregate kill
probability and, respectively, against $\delta_v$; the key inequality is
that $\beta$-kill implies $L$-extra-recovery because the $\beta$-kill
threshold is at most $\delta_v$ (attained when every in-neighbor is
active).

\begin{theorem}[Sandwich approximation, all $\beta$]
\label{thm:sandwich}
$L$ and $U$ are monotone submodular (Lemma~\ref{lem:recovery-submodular})
and $L\le f_\beta\le U$ pointwise (Lemmas~\ref{lem:upper}--\ref{lem:lower}).
Let $S_L, S_U$ be the greedy solutions for $L, U$ under $|S|\le k$, $S_\beta$
the greedy solution run directly on $f_\beta$, and
$S_{\mathrm{sand}}=\arg\max_{S\in\{S_L,S_U,S_\beta\}} f_\beta(S)$. Then, with
$S^\star$ optimal for $f_\beta$,
\begin{equation}
f_\beta(S_{\mathrm{sand}}) \ge \max\left\{\frac{f_\beta(S_U)}{U(S_U)},\ \frac{L(S^\star)}{f_\beta(S^\star)}\right\}\Big(1-\frac1e\Big) f_\beta(S^\star) \ge \rho\Big(1-\frac1e\Big)\mathrm{OPT}_\beta,
\end{equation}
where $\rho \eqdef \min_{|S|\le k} L(S)/U(S)$.
\end{theorem}
\emph{Proof.} Cited: the sandwich-approximation theorem of Lu, Chen, and
Lakshmanan~\cite{lu2015competition}, applied to the bracket $L\le f_\beta\le U$;
the $(1-1/e)$ factors are~\cite{nemhauser1978analysis} applied to greedy on
each submodular surrogate. The final inequality follows from
$L(S^\star)/f_\beta(S^\star)\ge L(S^\star)/U(S^\star)\ge\rho$. $\blacksquare$

Theorem~\ref{thm:sandwich} improves on treating backfire as an arbitrary
perturbation from submodularity via~\cite{horel2016maximization}, which
would only hold for $\beta=O(n^{-1/2})$ -- far below the
$\beta\in[0.1,0.7]$ regime our empirical study targets -- because it
exploits the \emph{structure} of the perturbation (a bounded extra kill
rate) rather than treating it adversarially, and it holds for every
$\beta$, recovering the tight $(1-1/e)$ exactly at $\beta=0$ (where $L=U$,
$\rho=1$). It also yields a practical corollary: since $L$ and $U$ are
ordinary $\beta=0$ processes that can be simulated directly, $f_\beta(S)/U(S)$
is a directly computable, per-instance certificate,
$f_\beta(S)\ge \frac{f_\beta(S)}{U(S)}(1-\frac1e)\mathrm{OPT}_\beta$ -- the
algorithm ships its own approximation guarantee on every instance it is run
on.

\begin{proposition}[Finite-horizon bound on $\rho$]
\label{prop:rho-bound}
$\rho \ge (1-\beta\Delta_{\mathrm{in}})^T$.
\end{proposition}
This is proven by charging every self-edge of a witness reachability path
under the coupling of Lemma~\ref{lem:lower}, and is conservative (it ignores
re-activation through alternate paths). A tighter steady-state conjecture,
$\rho \gtrsim q/(q+\beta\Delta_{\mathrm{in}})$ for $q=\min_v q_v$, is stated
in the underlying project notes as an explicit \emph{conjecture}: the
heuristic single-sojourn argument does not account for the downstream
cascade a shortened sojourn fails to seed, and closing that gap is left
open. This conjecture is not used in any proof in this paper (Section~\ref{sec:limitations}
discusses why it is, in any case, not the most promising target for
tightening the sandwich ratio).

\subsection{A closed-form index for the control problem}
\label{sec:index}

Given the mean-field belief field (Section~\ref{sec:algorithm}), each
node's per-round control problem is a 2-state MDP -- states
$s\in\{0,1\}$, discount $\gamma$, reward $w$ at $s=1$ -- with natural
transitions $a=p_{01}$ (activate), $b=p_{10}$ (deactivate), and three
actions: NONE (free, either state), CONVERT (cost $c_C$, forces $0\to1$,
only at $s=0$), MAINTAIN (cost $c_M$, forces stay-at-1, only at $s=1$),
under a shared Lagrange price $\lambda$ per unit cost. There are exactly
four deterministic stationary policies, named by (action at 0, action at
1): NN, NM, CN, CM.

\begin{proposition}[Closed forms]
\label{prop:closedform}
With $D\eqdef(1-\gamma)(1-\gamma+\gamma(a+b))$, each policy's value
functions $V_1(\lambda), V_0(\lambda)$ are given in closed, affine-in-$\lambda$
form by solving its fixed $2\times2$ linear system (e.g., for CN,
$V_0=-\lambda c_C+\gamma V_1$ and $V_1=w+\gamma[(1-b)V_1+bV_0]$, giving
$V_1=(w-\gamma b\lambda c_C)/((1-\gamma)(1+\gamma b))$); the full table of
four closed forms is given in the underlying project notes and reproduced
in the algorithm's implementation.
\end{proposition}

\begin{proposition}[Convexity in $\lambda$]
\label{prop:convexity}
$V^*(s,\lambda)=\max_\pi V_\pi(s,\lambda)$ over the four policies is a
maximum of four affine functions of $\lambda$: convex, piecewise linear,
with at most three breakpoints, computable in $O(1)$.
\end{proposition}
\emph{Proof.} In a discounted finite MDP some stationary deterministic
policy is optimal at every state simultaneously, so $V^*$ is exactly this
finite max; Killian, Perrault, and Tambe~\cite{killian2021beyond} (their
Prop.\ 4.1) prove convexity in $\lambda$ for the general multi-action
restless-bandit case, of which this is a one-line special case via finite
enumeration of the four policies. $\blacksquare$

\begin{proposition}[Index and indexability, this case]
\label{prop:indexability}
Define node $v$'s index $\lambda^\star_v$ as the largest $\lambda\ge0$ at
which the optimal policy still takes the paid action at $v$'s current state.
\end{proposition}
A proof \emph{sketch} shows the paid action's advantage over NONE has a
strictly negative slope in $\lambda$ from the cost term, while the value
gap $V^*_1-V^*_0$ is non-increasing in $\lambda$ on each linear piece,
checked by finite case analysis over the four policies. \textbf{We state
plainly, matching the underlying project notes' own status label, that
indexability -- i.e., that the paid set is a down-set $[0,\lambda^\star_v]$
-- is only verified numerically here} (on random instances, over a fine
$\lambda$ grid, checking that per-node cost is non-increasing in
$\lambda$), not proven in general; the finite four-policy structure makes
this tractable to check exactly where Killian et
al.~\cite{killian2021beyond} call general multi-action indexability
"notoriously difficult" and avoid it. The index itself is defined as the
supremum of the paid set, which remains well-defined and exactly computable
regardless.

\begin{theorem}[Per-round complexity]
\label{thm:complexity}
Given Propositions~\ref{prop:closedform}--\ref{prop:indexability}, one
round of the control policy costs $O(Km+n\log n)$, where $K$ is the number
of mean-field fixed-point iterations (geometric convergence under a
contraction condition on $\Delta\max(p^+,p^-)$).
\end{theorem}
The $n\log n$ term sorts nodes by $\lambda^\star_v$ and fills the shared
budget greedily -- exactly "smallest $\lambda$ with aggregate cost
$\le B$," the target of the bisection-plus-value-iteration procedure it
replaces, which cost $O(n\cdot I_{\mathrm{bisect}}\cdot I_{\mathrm{VI}})$
(on the order of $8000\,n$ scalar operations per round in our
configuration); this project's own measurements put the closed-form
replacement at roughly $130\times$ faster per round than that bisection
oracle. For comparison, naive Monte Carlo KKT greedy costs $O(knRm)$ per
seed set, and one static IMM selection costs
$O((k+\ell)(n+m)\log n/\varepsilon^2)$~\cite{tang2015influence}: our
per-round cost is near-linear like a single IMM run but is paid $T$ times,
the stated price of sustaining influence under decay rather than seeding
once.

\subsection{Hardness}
\label{sec:hardness}

For $\beta=0$ (and $q=0$, large $T$), the problem contains max-$k$-cover, so
$1-1/e$ is an unbeatable ceiling for every $\beta\ge0$ under
$\mathrm{P}\ne\mathrm{NP}$~\cite{feige1998threshold,nemhauser1978analysis} --
adding backfire cannot remove hard instances on which $\beta$ is simply
irrelevant. For adversarial, arbitrary non-submodular perturbations, no
constant-factor approximation exists in general once a function is
$\omega(n^{-1/2})$-far from submodular~\cite{horel2016maximization};
Theorem~\ref{thm:sandwich} escapes this only because our perturbation is
structured, not adversarial.

\begin{theorem}
\label{thm:hardness}
For every fixed $\beta\in[0,1)$, no polynomial-time algorithm approximates
$f_\beta$ within $(1-1/e+\varepsilon)$ for any $\varepsilon>0$ unless
$\mathrm{P}=\mathrm{NP}$.
\end{theorem}
\emph{Proof.} Reduction from max-$k$-cover,
$(1-1/e+\varepsilon)$-inapproximable unless
$\mathrm{P}=\mathrm{NP}$~\cite{feige1998threshold}. Build a bipartite
digraph (set-nodes to element-nodes, edge probability 1), each element
replicated to swamp the additive seed-count term. The key subtlety is that
backfire \emph{is} realized on this instance (a seed keeps attacking every
element it covers), so its harmlessness must be shown, not assumed. On this
instance class, an element with $c\ge1$ active covering seeds throughout is
a two-state chain with up-rate 1 and down-rate $1-(1-\beta)^c$, giving an
exact value formula
$f_\beta(S)=kT+T\sum_v\varphi_\beta(c_v(S))+O(n)$ with
$\varphi_\beta(c)=1/(2-(1-\beta)^c)$ strictly decreasing in $c\ge1$ and
bounded below by $1/2$; the max-coverage gap therefore survives backfire
intact. $\blacksquare$

This value formula also shows $f_\beta\ge f_0/2$ for \emph{every} seed set
on this instance class at $q=0$, exactly where the sandwich ratio $\rho$ of
Theorem~\ref{thm:sandwich} is vacuous ($\rho=0$).

\begin{proposition}[Conditional -- not a theorem]
\label{prop:conditional}
Suppose Feige's NO-instances additionally satisfy: for every candidate
$k$-subset $S$, the number of elements covered exactly once is at most
$(1/e+\varepsilon)n$. Then no polynomial-time algorithm achieves ratio
$c(\beta)(1-1/e)+\varepsilon'$, where
\begin{equation}
c(\beta)(1-1/e) = \frac1e + \Big(1-\frac2e\Big)\frac{\varphi_\beta(2)}{\varphi_\beta(1)}
= (1-1/e) - \Big(1-\frac2e\Big)\frac{\beta(1-\beta)}{1+2\beta-\beta^2}.
\end{equation}
\end{proposition}
Numerically $c(\beta)\approx0.968,0.946,0.939,0.940,0.960,0.981$ at
$\beta=0.1,0.25,0.4,0.5,0.75,0.9$, worst near $\beta\approx0.45$, vanishing
at both ends. \textbf{We flag this explicitly, and only, as a conditional
result, exactly as it is labeled in the underlying project notes}: the
extra property is not part of Feige's stated result, though it is
plausible (a Chernoff/union-bound concentration argument for the
"covered-exactly-once" statistic under a random partition), and two gaps
remain -- (i) the analogous block-averaging step needs concavity of
$d\mapsto\mathbb{E}\,\varphi_\beta(\mathrm{Bin}(d,1/k))$, unchecked; (ii)
cross-terms from two sets in the same partition class need the reduction's
pairwise-soundness property re-verified for this statistic. We report this
as "very likely true, not yet proven," and do not use it as a premise for
any theorem in this paper. Structural evidence -- a per-edge revive-at-least-as-fast-as-kill
asymmetry, and several natural gadget constructions all failing for the
same underlying reason -- suggests no super-constant worsening is
achievable for fixed $\beta$, but this too is offered as evidence, not a
proof.

\section{Algorithm: MF-BWI-Fair}
\label{sec:algorithm}

MF-BWI-Fair (Mean-Field Belief with a Lagrangian index and isoelastic-Welfare
reweighting, Fair) is a stateful per-round control policy. Each round:

\begin{enumerate}
\item \textbf{Bayesian posterior tracking.} A Beta--Bernoulli tracker
  maintains posterior means $\hat p_{uv}^+, \hat p_{uv}^-, \hat q_v$ for
  every edge/node from observed trial outcomes; the policy never sees the
  simulator's hidden true parameters, only these running estimates.
\item \textbf{Mean-field belief fixed point.} Exact per-node marginals
  require tracking correlations across the whole active-set distribution,
  intractable at scale. We use the standard individual-based / N-intertwined
  mean-field approximation from network epidemic models: each node's
  in-neighbors are replaced by their marginal activation probabilities
  $m_u$, giving the self-consistency equation
  \begin{equation}
  \begin{split}
  m_v = m_v(1-\hat q_v)\!\!\prod_{u\in N_{\mathrm{in}}(v)}\!\!(1-\hat p^-_{uv}m_u) \\
  {}+(1-m_v)\Big(1-\!\!\prod_{u\in N_{\mathrm{in}}(v)}\!\!(1-\hat p^+_{uv}m_u)\Big),
  \end{split}
  \end{equation}
  solved by Jacobi fixed-point iteration to tolerance $10^{-4}$ or 50
  iterations. We verified directly that the alternative network-coupling
  method of Ou et al.~\cite{ou2022networked} does not apply here, since it
  requires a deterministic action-vector coupling rather than the
  genuinely bilinear neighbor-\emph{state} coupling this model has.
\item \textbf{Group welfare weights.} Each group's current marginal
  welfare weight $w_g$ is computed from its tracked time-averaged realized
  reach $u_g$ (Section~\ref{sec:fairness-fix}), egalitarian by default.
\item \textbf{Closed-form Lagrangian index.} Using the converged belief
  field to obtain each node's "no-action" transition probabilities
  $p_{01}(v), p_{10}(v)$, each node's 2-state/3-action MDP is built with
  reward $r_v(s)=w_{g(v)}\cdot s$, and the shared-budget Lagrangian
  relaxation is solved via the closed-form index of
  Section~\ref{sec:index} (a bisection-plus-value-iteration oracle is kept
  as a test reference).
\item \textbf{Leftover-budget fill.} Any budget left over from the index
  scan's discreteness is filled by a greedy tie-break on each node's
  closed-form index (or Bellman gap, under the bisection oracle), keeping
  the hard budget constraint intact without reintroducing per-group floors
  or quotas.
\item \textbf{Execute, observe, update.} Actions are applied; the simulator
  advances one round under the layered-semantics transition rule
  (Section~\ref{sec:model}); observed trial outcomes update the Bayesian
  tracker; and each group's realized reach average $u_g$ is updated for the
  next round's weights.
\end{enumerate}

\begin{algorithm}[t]
\caption{MF-BWI-Fair, one round}
\label{alg:mfbwi}
\begin{algorithmic}[1]
\REQUIRE graph $G$, state $s$, budget $B$, fairness exponent $\alpha$
\STATE $(\hat p^+,\hat p^-,\hat q) \gets$ posterior means from tracker
\STATE $m \gets$ mean-field fixed point of Eq.~(5) from $(s,\hat p^+,\hat p^-,\hat q)$
\STATE $(p_{01}, p_{10}) \gets$ field transition probabilities from $m$
\STATE $u \gets$ tracked per-group time-averaged reach
\STATE $w \gets$ \textsc{GroupWelfareWeights}$(u, \alpha)$ \hfill (egalitarian default)
\FOR{each node $v$}
  \STATE build 2-state/3-action MDP with reward $w_{g(v)}\cdot s_v$, kernel $(p_{01}(v), p_{10}(v))$
  \STATE $\lambda^\star_v \gets$ closed-form index (Prop.~\ref{prop:indexability})
\ENDFOR
\STATE sort nodes by $\lambda^\star_v$ descending; select prefix fitting in $B$
\STATE greedily fill leftover budget by descending index among unselected nodes
\STATE apply CONVERT/MAINTAIN to selected nodes, NONE elsewhere
\STATE observe transition, update Bayesian tracker and $u_g$
\end{algorithmic}
\end{algorithm}

\subsection{The fairness fix: population-weighted vs.\ egalitarian welfare}
\label{sec:fairness-fix}

Rahmattalabi et al.'s isoelastic welfare weight, implemented exactly as
published, is $w_g=N_g u_g^{\alpha-1}$ (or $N_g/u_g$ at $\alpha=0$): the
correct marginal utility for maximizing \emph{aggregate} welfare, in which
a group's contribution scales with its population. This has a concrete,
measured downside for the stated goal of not neglecting a small group: at
$\alpha=0$, group $g$'s weight exceeds group $g'$'s only once
$u_{g'}/u_g > N_{g'}/N_g$ -- the smaller group must be under-served by a
factor proportional to the group-size \emph{ratio} before the mechanism
reacts at all. On this project's own three-group (20/40/60) benchmark
graph, at a neutral $\alpha_{\mathrm{fair}}=0$ this kept the smallest group
at a visibly lower realized reach than the largest purely because a
$3\times$ reach gap was required before the weight formula began favoring
it. We therefore adopt, as the default, the egalitarian
(population-unweighted) form $w_g=u_g^{\alpha-1}$ (or $1/u_g$ at
$\alpha=0$), a standard alternative welfare form in the social-welfare and
fair-division literature (not a new mechanism), counting each group once
regardless of size -- the appropriate choice when fairness means each
\emph{group}, not each person, is treated fairly regardless of population.

\section{Experimental Setup}
\label{sec:setup}

\textbf{Graph.} The primary validation graph is a single synthetic
stochastic block model with three groups of unequal size (20/40/60, 120
nodes total), within-block edge probability $p_{\mathrm{in}}=0.07$,
between-block edge probability $p_{\mathrm{out}}=0.008$, and per-edge
positive-influence probability $p_{uv}\sim\mathrm{Uniform}(0.05,0.15)$.
Density was hand-tuned so the network neither saturates near full
activation nor fails to cascade at all, either of which would mask the
effects under study. \textbf{We state plainly that this is the only graph
topology validated in this draft}; a multi-graph and real-dataset
validation is in progress (\texttt{experiments/MULTIGRAPH\_VALIDATION.md})
and is not reported here.

\textbf{Budget and horizon.} Per-round budget $B=100$, horizon $T=30$
rounds. One-shot baselines are given $k=20$ seeds, chosen so that
$k\cdot\mathrm{cost}(\mathrm{CONVERT})=100=B$: one round's worth of
MF-BWI-Fair's budget, spent once at round 0.

\textbf{Baselines.} \emph{KKT-greedy/CELF}~\cite{kempe2003maximizing,leskovec2007cost}:
classical Monte Carlo greedy under plain progressive IC ($\beta=0$, $q=0$),
this project's own reimplementation, used as the "ignore all four
relaxations" ablation. \emph{Fair-greedy}: an in-house welfare-weighted
greedy baseline in the style of~\cite{rahmattalabi2021fair} maximizing a
sum of concave per-group log-utilities. \emph{Repeated-greedy}: an in-house
cost-effective greedy that re-selects actions every round under the same
per-round budget, using \emph{true} parameters directly (no Bayesian
uncertainty, no fairness mechanism) -- the strongest plausible classical
competitor that also acts repeatedly, included specifically to isolate
whether MF-BWI-Fair's advantage is more than "gets to act every round."
\emph{IMM}~\cite{tang2015influence}: this project's independent
reimplementation of the two-phase martingale sampling and node-selection
algorithm, cross-checked against the paper's own pseudocode and against a
documented erratum~\cite{chen2018issue}. \emph{He--Kempe Saturate
Greedy}~\cite{he2016robust}: this project's independent reimplementation of
the bicriteria robust algorithm, given the low/high corners of the
$p_{uv}$ range as its scenario set. \emph{MF-BWI-Fair}, as described in
Section~\ref{sec:algorithm}.

\textbf{Honesty notes on baseline provenance.} Other than the objective
forms they target, every baseline in this comparison is an
independently-verified reimplementation, not the original authors' code.
Official code exists for IMM (a SourceForge project by the original
authors) but is not what is benchmarked here; no official code was found
for He--Kempe's Saturate Greedy or for Rahmattalabi et al.'s welfare
mechanism specifically. \textbf{A direct benchmark against IMM's original
C++ implementation is in progress}
(\texttt{experiments/IMM\_VALIDATION.md}) and is not reported here; until
it lands, all reported IMM numbers should be read as characterizing this
project's own reimplementation, cross-checked against the paper's stated
algorithm and theorem statements, not as reproducing the original authors'
published numbers.

\textbf{Metrics.} (a) time-averaged active-node count over the full
horizon; (b) per-group reach fraction (own-group-size-normalized,
time-averaged) and its minimum across groups; (c) wall-clock runtime.
Sweeps are conducted over backfire intensity $\beta$, recovery rate $q$,
and the fairness exponent $\alpha_{\mathrm{fair}}$, each with 15
independent trials using common random numbers across the swept parameter
within a trial to reduce sweep-trend noise without narrowing between-trial
variance.

\section{Results}
\label{sec:results}

% TODO: results section pending experiments/MULTIGRAPH_VALIDATION.md and experiments/IMM_VALIDATION.md

\section{Discussion and Limitations}
\label{sec:limitations}

\textbf{Single graph topology.} All findings to date are demonstrated on
one synthetic stochastic block model; generalization across graph
families, degree distributions, and real cascade datasets is the subject
of ongoing, not-yet-reported work.

\textbf{The conditional hardness result is exactly that.}
Proposition~\ref{prop:conditional} is not a theorem: it depends on one
plausible but unverified property of Feige's construction, and two
identified technical gaps (a concavity step and a cross-term
pairwise-soundness check) remain before it could become one. We report
the qualitative structural evidence -- the revive-at-least-as-fast-as-kill
asymmetry, and several gadget constructions all failing to encode a
hard-to-decide dependence on the seed set -- as suggestive that no
super-constant worsening exists for fixed $\beta$, not as a substitute for
a proof.

\textbf{No comparison against the original authors' code for He--Kempe or
Rahmattalabi et al.} Neither paper has published code corresponding to the
specific algorithm reimplemented here (a superficially related repository
for Rahmattalabi et al.\ in fact implements a different paper,
Tsang et al.~\cite{tsang2019group}, and must not be conflated with it).
Only IMM has verifiable original code, and even there this draft benchmarks
against our own reimplementation pending the validation noted in
Section~\ref{sec:setup}.

\textbf{The mean-field approximation's own error is not separately
quantified.} The belief fixed point of Section~\ref{sec:algorithm} is an
approximation whose deviation from the true (intractable) joint
activation distribution has not been bounded or measured in isolation from
the end-to-end policy's performance; this is a natural target for future
analysis.

\textbf{The sandwich bound is demonstrably loose, and this is a direction,
not a hidden flaw.} Exact computation on small instances (2--8 nodes,
$T\sim300$) shows the sandwich ratio $\rho$ loose by up to an order of
magnitude where it applies at all (e.g., roughly $13\times$ on a directed
star instance, because the lower surrogate inflates every node's recovery
uniformly by $\delta_v$ even where a node has no active in-neighbor left),
and vacuous whenever $q=0$ or $q+\beta\Delta_{\mathrm{in}}\ge1$. Separately,
$f_\beta/f_0$ is not bounded below by any instance-independent constant (a
bidirected $K_6$ instance at $\beta=0.5,q=0.05$ collapses by a factor of
roughly $300\times$ relative to $f_0$), so no purely multiplicative
comparison against $f_0$ can be the source of a strong general guarantee.
Yet greedy run on the plain $\beta=0$ surrogate landed within 1\% of
$\mathrm{OPT}_\beta$ on every one of the six small instances tested --
suggestive, not proof, that the practical quality of surrogate-greedy
substantially exceeds what $\rho(1-1/e)$ certifies. We read this as the
concrete next step for this line of work: a sharper \emph{upper}-bound
analysis (a per-instance or instance-class argument in the spirit of the
exact layered value formula behind Theorem~\ref{thm:hardness}), rather
than a search for a matching lower bound on $\rho$ itself, which the
evidence here suggests is not the right quantity to target.

\textbf{Indexability is checked, not proven.} As stated in
Proposition~\ref{prop:indexability}, the down-set property underlying the
closed-form index's exactness is verified numerically on random instances
over a fine $\lambda$ grid, not established as a general theorem for this
2-state/3-action case.

\textbf{An open, unconditional route to a matching lower bound.} The
underlying project notes identify, but do not pursue to completion, a
possible unconditional alternative to Proposition~\ref{prop:conditional}:
the class of functions $S\mapsto\sum_v\varphi(c_v(S))$ is closed under the
same symmetrization used in Vondr{\'a}k's symmetry-gap
framework~\cite{vondrak2013symmetry}, and the symmetric-instance gap for
$\varphi_\beta$ (Theorem~\ref{thm:hardness}'s proof) works out numerically
close to, and slightly below, Proposition~\ref{prop:conditional}'s
constant. Whether Vondr{\'a}k's smoothing-step argument -- which uses
submodularity directly -- extends to this non-submodular but structured
class is open and was not checked; if it does, it would give an
unconditional oracle-model lower bound without needing the unverified
property Proposition~\ref{prop:conditional} depends on. We flag this as the
single most promising next step toward turning Section~\ref{sec:hardness}'s
conditional result into a proof.

\section{Conclusion}

We formalized a combined relaxation of the four assumptions behind the
classical influence-maximization guarantee -- progressive dynamics,
monotone submodularity, known parameters, and fairness-free objectives --
and proved that a sandwich approximation recovers a $\rho(1-1/e)$ guarantee
for every backfire intensity, with a closed-form index making the
resulting control problem tractable at near-linear per-round cost. We
proved that $(1-1/e)$-hardness survives backfire exactly, while explicitly
declining to upgrade a plausible but unproven constant-factor worsening
result to theorem status. We corrected a measured fairness pathology in an
existing welfare objective and implemented the result in MF-BWI-Fair.
Preliminary single-graph experiments are encouraging; a full validation
against additional graph families and against IMM's original
implementation is in progress and left to future revisions of this
manuscript.

\bibliographystyle{IEEEtran}
\bibliography{references}

\end{document}
